\documentclass[11pt]{article}
\usepackage[margin=0.9in]{geometry}
\usepackage{graphicx,amsmath,amssymb,booktabs,hyperref,xcolor,tabularx}
\hypersetup{colorlinks=true,linkcolor=blue,urlcolor=blue}
\setlength{\parskip}{6pt}
\title{Method C v2: a Bayesian-symmetric monotone ansatz for the\\
       K=3 CRRA REE equilibrium price function\\
       (methodology + neural alternatives)}
\author{Fixed-Point Factory}
\date{June 2026}

\begin{document}
\maketitle

\section{The problem}

The K=3 CRRA Hellwig rational-expectations equilibrium is the
fixed-point equation
\[
P(u_1, u_2, u_3) \;=\; \Phi[P](u_1, u_2, u_3),
\]
where each agent observes a private signal $u_i$, market clearing
yields the price $P$, and $\Phi$ is a Bayesian-CRRA operator: it computes
each agent's posterior $\mu(p, u_k)$ from observing the price $p$ and their
own signal $u_k$, then CRRA-clears the market across the three agents.

\textbf{The unknowns} are the values of $P$. In our previous work we
represented $P$ as a 343-element tensor on a $G^3 = 7^3$ cube grid and
solved the FP equation cell-by-cell. We hit two structural problems:
\begin{itemize}
  \item The kernel-band Richardson operator (smooth, Newton-friendly) gives
        FPs that are \emph{non-monotone} in $\sim 27\%$ of cube cells at
        high $\tau$ --- non-physical artifacts.
  \item The strict-$h{=}0$ operator (Bayesianly correct) is non-smooth at
        $p$ values where the price level set passes through critical points
        of $P$; the FP iteration does not converge.
\end{itemize}

\section{The Bayesian symmetries}

The equilibrium price $P(u_1, u_2, u_3)$ must satisfy four structural
properties dictated by the model:

\begin{enumerate}
  \item \textbf{Permutation symmetry.} Agents are exchangeable:
        $P(u_1, u_2, u_3) = P(u_{\pi(1)}, u_{\pi(2)}, u_{\pi(3)})$ for any
        permutation $\pi$.
  \item \textbf{Sign-flip symmetry.} The state space is binary
        $v \in \{0, 1\}$ with symmetric prior; the signal distribution is
        $f_v(u) = N(\pm 0.5, 1/\tau)$, antisymmetric under $u \to -u$ and
        $v \leftrightarrow 1-v$. This gives
        $P(u_1, u_2, u_3) = 1 - P(-u_1, -u_2, -u_3)$.
  \item \textbf{Axis-wise monotonicity.} More positive $u_i$ increases the
        Bayesian posterior on $v=1$ for any agent observing it, and (via
        market clearing) raises the equilibrium price:
        $\partial P / \partial u_i > 0$ for each $i$.
  \item \textbf{Range.} $P \in (0, 1)$.
\end{enumerate}

Any candidate FP violating these is not a true equilibrium. The
cube-discretized solvers we used produce candidate FPs that
\emph{violate (3) and (4)} numerically. Instead of fighting the
discrete representation, Method C v2 builds (1)--(4) into the
parameterization.

\section{Method C v2 construction}

\textbf{Parameterize $P$ as}
\[
P(u_1, u_2, u_3) \;=\; \sigma\!\big( h(u_1) + h(u_2) + h(u_3) \big),
\qquad \sigma(x) = \frac{1}{1+e^{-x}},
\]
where $h: \mathbb{R} \to \mathbb{R}$ is a \textbf{1D odd, monotone-increasing}
function. Write
\[
h(u) \;=\; \int_0^u \sigma_*(s)^2 \, ds,
\]
with $\sigma_*: \mathbb{R} \to \mathbb{R}$ a 1D \emph{even} polynomial of
degree $2 d$:
\[
\sigma_*(s) \;=\; \sum_{n=0}^d a_n T_{2n}(s / U),
\]
in the Chebyshev basis (we use $U = 2.33$ to cover the practical signal
range).

\subsection{Why this enforces all four symmetries by construction}

\begin{enumerate}
  \item \textbf{Permutation symmetry.} $h(u_1) + h(u_2) + h(u_3)$ is
        manifestly symmetric in the agents.
  \item \textbf{Sign-flip symmetry.} $\sigma_*$ even $\Rightarrow \sigma_*^2$
        even $\Rightarrow h$ odd in $u$
        $\Rightarrow h(u_1) + h(u_2) + h(u_3)$ is odd in $(u_1, u_2, u_3) \to
        (-u_1, -u_2, -u_3)$
        $\Rightarrow$ $\sigma(L) + \sigma(-L) = 1$.
  \item \textbf{Monotonicity.} $\sigma_*^2 \geq 0$ gives $h'(u) = \sigma_*^2(u) \geq 0$,
        so $h$ is monotone non-decreasing. Hence
        $\partial P / \partial u_i = \sigma'(L) \cdot h'(u_i) \geq 0$.
  \item \textbf{Range.} $\sigma \in (0, 1)$ for any finite $L$.
\end{enumerate}

\subsection{Unknowns}

Just $d+1$ scalars: $a_0, a_1, \dots, a_d$. We use $d = 4$ (so $\sigma_*$ is
degree 8) giving \textbf{5 unknowns per cell}. Vastly fewer than the
343 cube unknowns or the 84 symmetric-reduced cube unknowns.

\subsection{Numerical pipeline per $(\gamma, \tau)$ cell}

\begin{enumerate}
  \item Given current $\{a_n\}$:
  \begin{itemize}
    \item Compute $h(u_i)$ at the $G = 7$ cube nodes via 16-point
          Gauss-Legendre quadrature on $[0, u_i]$.
    \item Form $P_{ijk} = \sigma(h(u_i) + h(u_j) + h(u_k))$.
  \end{itemize}
  \item Apply the standard Lin-CDF R4 kernel-band operator $\Phi$ to get
        $\Phi[P]$.
  \item Compute residual $r_{ijk} = \Phi[P]_{ijk} - P_{ijk}$.
  \item Minimize $\|r\|_2^2$ over the 5 unknowns via SciPy
        \texttt{least\_squares(method=trf, jac="2-point")}.
\end{enumerate}

\textbf{Per cell wall time}: $\sim 1$ s (5 unknowns + few-point FD Jacobian
+ a handful of TRF iterations).

\subsection{What the residual means}

The minimised residual $|F|_\infty = \max_{ijk} |\Phi[P]_{ijk} - P_{ijk}|$
is the \textbf{structural non-additivity} of the equilibrium: how much it
deviates from being expressible as a sum of three identical 1D
contributions. This separates cleanly from:
\begin{itemize}
  \item Kernel-band bias (gone --- since we measure the residual at the
        \emph{ansatz} $P$, not at a kernel-band-biased FP).
  \item Discretisation noise (gone --- the ansatz is analytic).
  \item Operator non-smoothness (gone --- our $P$ is $C^\infty$).
\end{itemize}

If $F \to 0$ at low $\tau$, the equilibrium IS additive --- can be written
as a sum of agent-wise contributions. If $F$ grows at high $\tau$, the
equilibrium has irreducible joint structure (genuine cross-signal
information beyond the sum of marginal posteriors).

\section{Empirical results (preliminary)}

From a partial $\tau$ ladder run at $\gamma = 100$:
\begin{center}
\begin{tabular}{rrll}
\toprule
$\tau$ & residual $|F|_\infty$ & meaning & wall \\
\midrule
0.025 & $2.7 \times 10^{-8}$ & ansatz captures equilibrium nearly exactly & 1 s \\
0.05  & $2.1 \times 10^{-7}$ & still additive to high precision & 1 s \\
0.10  & $5.5 \times 10^{-6}$ & beginning to see non-additivity & 1 s \\
0.25  & $8.0 \times 10^{-4}$ & moderate non-additivity & 1 s \\
0.50  & $3.4 \times 10^{-2}$ & substantial non-additivity & 1 s \\
1.00  & $\sim 5 \times 10^{-2}$ (proj.) & strong non-additivity, P spans $\sim$ full range & 1 s \\
\bottomrule
\end{tabular}
\end{center}

The residual rises smoothly with $\tau$, suggesting the equilibrium is
\textbf{additive to high precision at low $\tau$} and acquires irreducible
joint structure as the signal precision grows.

\section{Could a neural network do better?}

\textbf{Universal approximation: yes.} A neural network is a universal
function approximator, so in principle it could fit the equilibrium to
arbitrary precision. But the right question is: what's the
\emph{efficient} parameterisation that respects the same Bayesian
symmetries?

\subsection{Naïve MLP: no}

A vanilla MLP $\mathrm{NN}(u_1, u_2, u_3)$ has none of the four required
symmetries built in. Enforcing them via training:
\begin{itemize}
  \item Permutation symmetry: data augmentation, but expensive and only
        approximate;
  \item Sign-flip: similar issues;
  \item Monotonicity: requires constrained training (e.g., projected
        gradient or monotone-NN architectures with positive weights and
        monotone activations like ReLU/softplus everywhere).
\end{itemize}
This is much harder than our 5-parameter polynomial.

\subsection{Symmetric monotone NN: yes, but redundant}

The natural NN counterpart of Method C v2 is
\[
P(u_1, u_2, u_3) = \sigma\!\big( \phi(u_1) + \phi(u_2) + \phi(u_3) \big),
\]
where $\phi: \mathbb{R} \to \mathbb{R}$ is a 1D \emph{monotone}
neural net (e.g., a stack of positive-weight linear layers with
positive activations, or a monotone PCHIP NN). This automatically gives
the same four symmetries.

For our problem, the empirical evidence is that a degree-8 even polynomial
(Cheby basis, 5 coefficients) already captures $\phi$ to F $\sim 10^{-6}$
at low $\tau$. An NN would add hundreds of parameters with diminishing
returns. \textbf{The 1D polynomial is a near-optimal basis here.}

\subsection{Joint correction: yes, useful but constrained}

To go beyond the additive ansatz at high $\tau$, add a joint correction:
\[
P = \sigma\!\big( h(u_1) + h(u_2) + h(u_3) + g(u_1, u_2, u_3) \big),
\]
where $g$ is a permutation-symmetric, sign-flip-symmetric function
(e.g., $g$ symmetric polynomial in $u_i$, odd in the overall sign flip).
For $g$, NN choices:
\begin{itemize}
  \item \textbf{Deep Sets} architecture: $g = \psi(\sum_i \rho(u_i))$
        with $\rho, \psi$ small NNs. Gives permutation invariance and
        moderate flexibility. ~$10^2$ params.
  \item \textbf{Equivariant NN}: stricter constraint than Deep Sets, fewer
        params, faster.
  \item \textbf{Symmetric polynomial basis}: $g = \sum c_{\mathbf{m}}
        e_{\mathbf{m}}(u_1, u_2, u_3)$ in elementary symmetric polynomials,
        with sign-flip-respecting parity. ~$10$--$50$ params.
\end{itemize}

We expect the symmetric polynomial basis to be most parsimonious for
the K=3 problem and to match NN performance with fewer parameters.

\subsection{Bottom line on NNs}

\begin{itemize}
  \item NNs are \emph{not} a clear win over the polynomial-Cheby ansatz
        for this problem. The symmetries are exact (not just learnable)
        with the 5-parameter polynomial.
  \item For the \emph{joint correction} to capture remaining
        non-additivity at high $\tau$, both NN and symmetric polynomial
        bases would work. Polynomial wins on parsimony for the symmetric
        structure we know holds.
  \item Worth trying only if the polynomial bases plateau at a residual
        we want to push lower.
\end{itemize}

\section{Reproducibility}

\begin{itemize}
  \item Operator + ansatz:
        \texttt{projects/REZN/solver\_code/highprec\_dd\_qd/dd\_k3\_method\_c\_v2.py}
  \item Ladder driver:
        \texttt{projects/REZN/solver\_code/highprec\_dd\_qd/dd\_k3\_method\_c\_v2\_ladder.py}
  \item Per-$\tau$ saved coefficients:
        \texttt{projects/REZN/solved\_fixed\_points/dd\_k3\_overnight/method\_c\_v2/coefs/tau\{tau\}.npy}
\end{itemize}

\end{document}
